% Source: https://artofproblemsolving.com/community/c3h2313126_hexagon_and_dodecagon_in_a_circle
% Stable identity: aops-topic-2313126-q01
Let the circle have centre $O$ and radius $R$.

\textbf{The hexagon.} Joining $O$ to the six vertices cuts the hexagon into six equilateral triangles of side $R$, each of area $\dfrac{\sqrt3}{4}R^2$. So
\[\dfrac{3\sqrt3}{2}R^2=18\sqrt3,\qquad R^2=12.\]

\textbf{The dodecagon.} Joining $O$ to the twelve vertices cuts it into twelve congruent isosceles triangles such as $OAB$ (shaded), with $OA=OB=R$ and $\angle AOB=\dfrac{360^\circ}{12}=30^\circ$.

\begin{asy}
usepackage("amsmath");
size(6cm);
defaultpen(fontsize(10pt));
real R=1;
draw(circle((0,0),R),gray);
pair[] H; for(int k=0;k<6;++k) H.push(R*dir(60k));
pair[] D; for(int k=0;k<12;++k) D.push(R*dir(30k));
draw(H[0]--H[1]--H[2]--H[3]--H[4]--H[5]--cycle,blue+linewidth(0.9));
draw(D[0]--D[1]--D[2]--D[3]--D[4]--D[5]--D[6]--D[7]--D[8]--D[9]--D[10]--D[11]--cycle,red);
pair O=(0,0), A=D[1], B=D[0], F=(A.x,0);
fill(O--B--A--cycle,palegreen);
draw(O--B--A--cycle,heavygreen+linewidth(0.9));
draw(A--F,dashed+heavygreen);
draw(F+(0,0.05)--F+(-0.05,0.05)--F+(-0.05,0),heavygreen);
dot(O); label("$O$",O,W); label("$A$",A,NE); label("$B$",B,E); label("$F$",F,S,fontsize(8pt));
draw(arc(O,0.22,0,30)); label("$30^\circ$",0.24*dir(15),E,fontsize(7pt));
label("$R$",(O+A)/2,NW,fontsize(8pt));
\end{asy}

Drop the perpendicular $AF$ from $A$ to $OB$. Triangle $OFA$ is a $30^\circ$-$60^\circ$-$90^\circ$ triangle with hypotenuse $OA=R$, so $AF=\dfrac R2$. Hence
\[[OAB]=\dfrac12\cdot OB\cdot AF=\dfrac12\cdot R\cdot\dfrac R2=\dfrac{R^2}{4},\]
and the dodecagon has area
\[12\cdot\dfrac{R^2}{4}=3R^2=3\cdot12=\boxed{36}\text{ cm}^2.\]
